\documentclass[a4paper,10pt]{article}
\usepackage[francais]{babel}
\usepackage{url}

\input{letter-icube-uds.sty}

\enCouleur

\begin{document}

\name{Stéphane \textsc{Genaud}}
\fonction{Professeur ENSIIE, mise à disposition Université de Strasbourg}
\telephoneillkirch{+33 (0)3688 54542}
\email{genaud@unistra.fr}

\begin{letter}[Recommandation de Tayeb Nazim DJERADI]{à}{Fondation Francis Bouygues}

Madame, Monsieur,

  J'ai été sollicité Tayeb Nazim DJERADI pour cette lettre de recommandation.

  Tayeb Nazim DJERADI est étudiant en première année de Licence Informatique,
  à la faculté de Mathématiques et Informatique de l'université de Strasbourg.
  
  Je connais cet étudiant depuis la rentrée universitaire 2020-2021 (septembre
  2020) en tant qu'enseignant. J'ai eu le plaisir de l'avoir dans mon groupe
	  d'étudiants dans le cours \emph{Algorithmique et Programmation}.

  Je suis également responsable de ce cours, et bien que je ne puisse vous communiquer
  officiellement sa note avant le jury, je peux dire factuellement que ses résultats
  scolaires dans cette matière, importante dans le cursus, sont excellents, puisqu'il
	  se classe 19e sur les 324 étudiants qui ont eu une note dans cette matière 
	  (note sur trois épreuves).

  Ce résultat me conforte dans la très bonne impression que Nazim m'a fait durant
  nos séances de cours en présentiel : il a fait preuve d'une assiduité irréprochable, 
  d'une grande curiosité, et s'est révélé être un moteur pour le groupe.
  Je suis persuadé que Nazim a la motivation et la capacité à réaliser un
  excellent parcours scolaire qui pourra, je l'espère, aller jusqu'à un Master.

	  Pour ces raisons, je recommande chaleureusement sa candidature pour une bourse 
	  de votre fondation.

   Veuillez agrééer, Madame, Monsieur, l'expression de mes sentiments respectueux.

\begin{flushright}
\includegraphics[width=.25\textwidth]{signgenaud.jpg} 
\end{flushright}

\end{letter}
}

%\newpage
%\bibliographystyle{plain}
%\bibliography{arno}


\end{document}
